\documentclass[11pt,a4paper]{article}

\usepackage[margin=2.5cm]{geometry}
\usepackage{graphicx}
\usepackage{hyperref}
\usepackage{booktabs}
\usepackage{longtable}
\usepackage{listings}
\usepackage{xcolor}
\usepackage{parskip}

\hypersetup{
    colorlinks=true,
    linkcolor=blue,
    urlcolor=blue
}

\lstset{
    basicstyle=\ttfamily\small,
    breaklines=true,
    frame=single,
    backgroundcolor=\color{gray!10},
    columns=fullflexible
}

\title{\includegraphics[width=3cm,height=3cm,keepaspectratio]{image.png}\\[1em] Cryptography and Network Security \\ \large Integrated Situation - Technical Report}
\author{4202670024 GBOUSSOUBO IGNAZIANKI \\ ETT B-Tech, Year 4, ULK Polytechnic Institute}
\date{\today}

\begin{document}

\maketitle

\begin{abstract}
This report documents the security assessment and toolkit developed for the
ULK Polytechnic student records system, covering risk assessment, an
encryption and integrity-checking programme, and network traffic filtering
rules.
\end{abstract}

\tableofcontents
\newpage
\section{Risk Assessment}

\subsection{Assets, Vulnerabilities, and Consequences}

\begin{table}[h]
\centering
\begin{tabular}{p{3.5cm}p{4cm}p{6cm}}
\toprule
\textbf{Asset} & \textbf{Vulnerability} & \textbf{Consequence} \\
\midrule
Student records database & Weak staff passwords & Attacker guesses/brute-forces login, leading to data breach or record tampering \\
Inter-campus file transfer & Unencrypted transfer & Data intercepted or altered in transit (man-in-the-middle) \\
Records server & Outdated software combined with guest network access & Unpatched server reachable by anyone on the guest network, leading to full compromise \\
\bottomrule
\end{tabular}
\caption{Assets, vulnerabilities, and consequences}
\end{table}

\subsection{Risk Ranking}

\begin{table}[h]
\centering
\begin{tabular}{p{0.7cm}p{4cm}p{1.8cm}p{1.8cm}p{5cm}}
\toprule
\textbf{Rank} & \textbf{Risk} & \textbf{Likelihood} & \textbf{Impact} & \textbf{Reason} \\
\midrule
1 & Outdated software + guest access & High & High & No credentials needed; matches the suspicious external connection attempts already observed \\
2 & Weak staff passwords & Medium-High & High & Easy to exploit, gives full legitimate-looking access \\
3 & Unencrypted file transfer & Medium & Medium--High & Attacker needs network-path position first \\
\bottomrule
\end{tabular}
\caption{Risk ranking by likelihood and impact}
\end{table}

\subsection{Recommended Controls}

\begin{table}[h]
\centering
\begin{tabular}{p{5cm}p{8cm}}
\toprule
\textbf{Risk} & \textbf{Control} \\
\midrule
Outdated software + guest access & Segment guest network away from the server; apply regular patching \\
Weak passwords & Enforce a strong password policy and multi-factor authentication (MFA) \\
Unencrypted transfer & Encrypt transfers (SFTP/TLS) and encrypt files before sending \\
\bottomrule
\end{tabular}
\caption{Recommended controls}
\end{table}

\section{Encryption and Integrity Toolkit}

A Python command-line toolkit, \texttt{security\_toolkit.py}, was developed
using the \texttt{cryptography} library (Fernet symmetric encryption) and
the built-in \texttt{hashlib} module (SHA-256).

\subsection{Functions}

\begin{itemize}
    \item \textbf{Key generation}: creates a random Fernet key, stored outside the repository (\texttt{key.key}, excluded via \texttt{.gitignore}).
    \item \textbf{Encryption}: reads a student record file and writes an encrypted \texttt{.enc} version using the key.
    \item \textbf{Decryption}: reverses the process, reconstructing the exact original file from the \texttt{.enc} file and the key.
    \item \textbf{Integrity check}: computes a SHA-256 hash of a file and can later re-check the file against the saved hash to detect any modification.
    \item \textbf{Error handling}: every operation checks for missing files, missing keys, and invalid/wrong keys, and exits with a clear error message instead of crashing.
\end{itemize}

\subsection{Test Evidence}

The toolkit was run end-to-end against a sample student record file
(\texttt{sample\_student\_record.csv}):

\begin{lstlisting}
$ python3 security_toolkit.py encrypt sample_student_record.csv sample_student_record.enc --key-file key.key
Encrypted 'sample_student_record.csv' -> 'sample_student_record.enc'

$ python3 security_toolkit.py decrypt sample_student_record.enc sample_student_record.decrypted.csv --key-file key.key
Decrypted 'sample_student_record.enc' -> 'sample_student_record.decrypted.csv'

$ diff sample_student_record.csv sample_student_record.decrypted.csv
(no output -- files are identical)

$ python3 security_toolkit.py hash sample_student_record.csv --save sample_student_record.sha256
SHA-256 of 'sample_student_record.csv': 55e91847c70d981356b9c6f522e60c53c8e86e409598886389fad893f4dfd609
Saved to 'sample_student_record.sha256'

$ python3 security_toolkit.py verify sample_student_record.csv --hash-file sample_student_record.sha256
OK: 'sample_student_record.csv' matches the recorded hash. No tampering detected.
\end{lstlisting}

This confirms: (1) encryption and decryption are correct and lossless, since
the decrypted file is byte-for-byte identical to the original, and (2) the
SHA-256 integrity check correctly identifies an unmodified file, and (as
tested separately) correctly flags a modified one with a mismatch warning.

\section{Network Traffic Filtering}

\subsection{Environment Note}

No physical or virtual lab network was available for this exam. The rules
below are the actual \texttt{iptables} commands that would be applied on
the records server (or its gateway) in a real deployment. To still produce
genuine, reproducible test evidence, a Python simulator
(\texttt{firewall\_simulator.py}) re-implements the identical rule logic,
evaluated in the same top-to-bottom order as \texttt{iptables}, and runs
real connection-attempt tests against it.

\subsection{Network Assumptions}

\begin{table}[h]
\centering
\begin{tabular}{p{6cm}p{6cm}}
\toprule
\textbf{Item} & \textbf{Value} \\
\midrule
Records server & 10.0.0.10 \\
Service port & 8443 \\
Staff network (authorised) & 10.0.1.0/24 \\
Guest network (blocked) & 10.0.2.0/24 \\
Suspicious external address & 203.0.113.50 \\
\bottomrule
\end{tabular}
\caption{Network assumptions used for the firewall configuration}
\end{table}

\subsection{Firewall Rules}

\begin{lstlisting}
# (a) Block guest network access to the student records server
iptables -A INPUT -s 10.0.2.0/24 -d 10.0.0.10 -j DROP

# (b) Permit authorised staff network access to the service
iptables -A INPUT -s 10.0.1.0/24 -d 10.0.0.10 -p tcp --dport 8443 -j ACCEPT

# (c) Block other inbound access to that service
iptables -A INPUT -d 10.0.0.10 -p tcp --dport 8443 -j DROP

# Default policy
iptables -P INPUT DROP
\end{lstlisting}

\subsection{Test Results}

\begin{table}[h]
\centering
\begin{tabular}{p{0.5cm}p{5cm}p{2cm}p{2cm}p{1.5cm}}
\toprule
\textbf{\#} & \textbf{Test} & \textbf{Expected} & \textbf{Actual} & \textbf{Result} \\
\midrule
1 & Staff $\rightarrow$ records service (10.0.1.50) & ACCEPT & ACCEPT & PASS \\
2 & Guest $\rightarrow$ records service (10.0.2.75) & DROP & DROP & PASS \\
3 & Unknown external host $\rightarrow$ service (203.0.113.50) & DROP & DROP & PASS \\
\bottomrule
\end{tabular}
\caption{Firewall test results (simulated)}
\end{table}

All three tests passed, confirming the rules correctly permit staff traffic
to the service while blocking guest and unrecognised external traffic.

\section{Conclusion}

This project identified and ranked three key risks to the student records
system -weak staff passwords, unencrypted inter-campus file transfer, and
an outdated, guest-accessible server  and proposed practical controls for
each. A Python toolkit was implemented and tested to provide file
encryption, decryption, and SHA-256 integrity verification, successfully
demonstrating correct round-trip encryption and tamper detection. Firewall
rules were designed to restrict access to the records server to authorised
staff only, and validated through a rule-equivalent simulation, with all
three required tests passing.

\vspace{1em}
\noindent\textbf{GitHub repository:} \url{https://github.com/Itzspencer235/cryptography-network-security-exam}

\end{document}